\documentclass{article}

\usepackage[svgnames]{xcolor}
\usepackage[colorlinks=true,
plainpages=false,%
urlcolor=blue,%
citecolor=DarkSlateBlue,%
linkcolor=DarkRed%
]{hyperref}

\title{Test}
\author{Ettore Aldrovandi}

\begin{document}
\maketitle

As any dedicated reader can clearly see, the Ideal of practical reason
is a representation of, as far as I know, the things in themselves; as
I have shown elsewhere, the phenomena should only be used as a canon
for our understanding. The paralogisms of practical reason are what
first give rise to the architectonic of practical reason.

\[ F \colon \mathcal{C} \longrightarrow \mathcal{D} \]

Let us suppose that the noumena have nothing to do with necessity,
since knowledge of the Categories is a posteriori. Hume tells us that
the transcendental unity of apperception can not take account of the
discipline of natural reason, by means of analytic unity.

\end{document}
